% \documentclass[bachelor,assignment]{pmk-graduate}
\documentclass[bachelor,paper]{pmk-graduate}

%%%%%%%%%%%%%%%%%%%%%%%%%%%%%%%%%%%%%%%%%%%%%%%%%%%%%%%%%%%%%%%%%
% Подключаете и определяете в соответствии с вашими потребностями
%%%%%%%%%%%%%%%%%%%%%%%%%%%%%%%%%%%%%%%%%%%%%%%%%%%%%%%%%%%%%%%%%
\usepackage{amsthm}
\usepackage{amsmath}
\usepackage[shortcuts]{extdash}
\usepackage{hyphenat}
\usepackage{graphicx}
\usepackage{chngcntr}
\usepackage{float}
% \usepackage{mathtext}
\usepackage{amssymb}
% \usepackage{algorithm,algorithmicx,algpseudocode}
\usepackage{algorithmicx}

\newtheorem{thm}{Теорема}
\theoremstyle{definition}
\newtheorem{dfn}{Определение}
%%%%%%%%%%%%%%%%%%%%%%%%%%%%%%%%%%%%%%%%%%%%%%%%%%%%%%%%%%%%%%%%%

%%%%%%%%%%%%%%%%%%%%%%%%%%%%%%%%%%%%%%%%%%%%%%%%%%%%%%%%%%%%%%%%%
% Список литературы ВКР
%%%%%%%%%%%%%%%%%%%%%%%%%%%%%%%%%%%%%%%%%%%%%%%%%%%%%%%%%%%%%%%%%
\addbibresource{references.bib}
%%%%%%%%%%%%%%%%%%%%%%%%%%%%%%%%%%%%%%%%%%%%%%%%%%%%%%%%%%%%%%%%%

\begin{document}
% Учебный год
\year{2025-2026}

% Приказ: дата, номер
%\order{02.12.2025}{222-С}

% Студент
\student{Щербаков Матвей Андреевич
  \brief М.\,А.\,Щербаков
  \group 47
  \english Scherbakov Matvey Andreevich
}

% Научный руководитель
\superviser{Солдатенко Илья Сергеевич
  \brief И.\,С.\,Солдатенко
  \position доцент кафедры информационных технологий
  \degree к.\,ф.-м.\,н. % Учёная степень. При отсутствии убрать строку
  \rank доцент % Учёное звание. При отсутствии убрать строку
}

% Образовательная программа. См. перечень ниже
\program{swi}
% Бакалавриат:
% ics - Информатика и компьютерные науки
% swi - Инженерия программного обеспечения
% swiai - Программная инженерия в искусственном интеллекте
% aida - Искусственный интеллект и анализ данных
% mm - Математическое моделирование
% sa - Системный анализ
% icpim - Инженерное компьютерное проектирование и индустриальная математика
% acse - Прикладная информатика в экономике
% acsm - Прикладная информатика в мехатронике
% icmrs - Интеллектуальное управление в мехатронных и робототехнических системах

% Заголовок ВКР
\title{Мобильное приложение для мониторинга городских проблем на основе технологии React Native
\english Mobile application for monitoring urban problems using React Native technology}

% Цель ВКР
\goal{Разработать мобильное приложение на основе технологии React Native для мониторинга и отслеживания городских проблем с возможностью создания заявок, просмотра объявлений.}

% Задачи для разработки в ВКР
\begin{tasks}

  \task{Анализ предметной области и существующих решений.
    \result
    Проведен анализ существующих платформ для мониторинга городских проблем. Изучены аналогичные решения, их функциональные возможности и недостатки. Определены требования к разрабатываемому приложению.
  }

  \task{Выбор и обоснование технологий разработки.
    \result
    Произведен сравнительный анализ технологий кроссплатформенной разработки. Обоснован выбор React Native в качестве основного фреймворка разработки. Определен технологический стек проекта.
  }

  \task{Проектирование архитектуры мобильного приложения.
    \result
    Разработана архитектура приложения, определены основные компоненты и модули системы. Спроектирована структура данных и взаимодействие с backend API. Разработана навигационная структура приложения.
  }

  \task{Разработка пользовательского интерфейса приложения.
    \result
    Спроектированы и реализованы экраны приложения: раздел с объявлениями, интерактивная карта с отметками проблем, список пользовательских заявок, технический раздел.
  }

  \task{Реализация функционала работы с картой и кластеризации маркеров.
    \result
    Интегрирован картографическая библиотека в приложение. Реализована система отображения маркеров на карте с использованием библиотеки кластеризации для оптимизации производительности при большом количестве заявок.
  }

  \task{Реализация системы создания и управления заявками.
    \result
    Разработан функционал создания новых заявок с возможностью выбора местоположения на карте, добавления описания, фотографий и категории проблемы. Реализован раздел управления пользовательскими заявками с отображением статуса и обратной связи от администрации.
  }

\end{tasks}

% Перечень иллюстративного материала
%\begin{images}
%\image{1}
%\image{2}
%\end{images}

% Консультанты
%\begin{advisers}
%\adviser{Сидоров Сидор Сидорович, проектирование схемы БД}
%\end{advisers}

% Команда для генерации задания на ВКР (закомментируйте, если нужно скомпилировать текст ВКР)
% \makeassignment

\text{привет мир}

\end{document}
